\documentclass[11pt,a4paper]{article}
\usepackage[top=2.3cm,bottom=2.3cm,left=2.2cm,right=2.2cm]{geometry}

\usepackage{amsmath,amssymb}
\usepackage{fontspec}
\setmainfont{Liberation Serif}
\setsansfont{Liberation Sans}
\setmonofont{DejaVu Sans Mono}[Scale=0.84]

\usepackage{xcolor}
\definecolor{brandblue}{RGB}{20, 55, 105}
\definecolor{accentblue}{RGB}{35, 95, 165}
\definecolor{codebg}{RGB}{248, 249, 251}
\definecolor{codeframe}{RGB}{215, 222, 230}
\definecolor{javakeyword}{RGB}{0, 45, 170}
\definecolor{javastring}{RGB}{5, 120, 20}
\definecolor{javacomment}{RGB}{120, 120, 120}
\definecolor{javanumber}{RGB}{20, 75, 220}

\definecolor{noteborder}{RGB}{30, 110, 65}
\definecolor{notebg}{RGB}{242, 249, 245}
\definecolor{warnborder}{RGB}{185, 75, 10}
\definecolor{warnbg}{RGB}{254, 248, 240}
\definecolor{tipborder}{RGB}{30, 95, 160}
\definecolor{tipbg}{RGB}{242, 247, 254}

\usepackage{fancyhdr}
\usepackage{lastpage}
\pagestyle{fancy}
\fancyhf{}
\lhead{\parbox[t]{0.58\textwidth}{\raggedright\small\sffamily\color{brandblue}\textbf{T08 --- Costruttori, Inizializzazione ed Incapsulamento}}}
\rhead{\small\sffamily\color{gray}Programmazione II -- UniVR (2025/2026)}
\cfoot{\small\sffamily Pagina \thepage\ di \pageref{LastPage}}
\renewcommand{\headrulewidth}{0.5pt}
\renewcommand{\headrule}{\hbox to\headwidth{\color{brandblue!70!black}\leaders\hrule height \headrulewidth\hfill}}

\usepackage{titlesec}
\titleformat{\section}{\Large\bfseries\sffamily\color{brandblue}}{\thesection}{0.8em}{}[\titlerule]
\titleformat{\subsection}{\large\bfseries\sffamily\color{accentblue}}{\thesubsection}{0.8em}{}
\titleformat{\subsubsection}{\normalsize\bfseries\sffamily\color{brandblue!85!black}}{\thesubsubsection}{0.8em}{}

\usepackage{needspace}
\usepackage{fancyvrb}
\DefineVerbatimEnvironment{asciibox}{Verbatim}{
  fontsize=\fontsize{7.5pt}{9.2pt}\selectfont,
  frame=single,
  rulecolor=\color{codeframe},
  fillcolor=\color{codebg},
  framesep=2.5mm
}
\usepackage{listings}
\lstset{
  backgroundcolor=\color{codebg},
  frame=single,
  rulecolor=\color{codeframe},
  basicstyle=\ttfamily\footnotesize,
  keywordstyle=\color{javakeyword}\bfseries,
  stringstyle=\color{javastring},
  commentstyle=\color{javacomment}\itshape,
  numberstyle=\tiny\color{gray},
  numbers=left,
  numbersep=7pt,
  stepnumber=1,
  breaklines=true,
  breakatwhitespace=true,
  showstringspaces=false,
  tabsize=4,
  xleftmargin=12pt,
  framexleftmargin=12pt
}

\newcommand{\passthrough}[1]{#1}
\providecommand{\tightlist}{\setlength{\itemsep}{0pt}\setlength{\parskip}{0pt}}

\usepackage{tcolorbox}
\tcbuselibrary{skins,breakable}
\newtcolorbox{studynote}[1][Nota]{
  enhanced,
  breakable,
  colback=notebg,
  colframe=noteborder,
  fonttitle=\bfseries\sffamily\small,
  coltitle=white,
  title={#1},
  arc=2mm,
  boxrule=0.8pt,
  left=9pt, right=9pt, top=7pt, bottom=7pt
}

\newtcolorbox{studywarning}[1][Attenzione]{
  enhanced,
  breakable,
  colback=warnbg,
  colframe=warnborder,
  fonttitle=\bfseries\sffamily\small,
  coltitle=white,
  title={#1},
  arc=2mm,
  boxrule=0.8pt,
  left=9pt, right=9pt, top=7pt, bottom=7pt
}

\newtcolorbox{studytip}[1][Suggerimento]{
  enhanced,
  breakable,
  colback=tipbg,
  colframe=tipborder,
  fonttitle=\bfseries\sffamily\small,
  coltitle=white,
  title={#1},
  arc=2mm,
  boxrule=0.8pt,
  left=9pt, right=9pt, top=7pt, bottom=7pt
}

\usepackage{booktabs}
\usepackage{longtable}
\usepackage{array}
\usepackage{calc}

\usepackage{hyperref}
\hypersetup{
  colorlinks=true,
  linkcolor=accentblue,
  urlcolor=accentblue,
  citecolor=brandblue,
  pdfborder={0 0 0}
}

\title{\Huge\bfseries\sffamily\color{brandblue} T08 --- Costruttori, Inizializzazione ed Incapsulamento \\[0.3em] \large\sffamily Ciclo di Inizializzazione, Overloading dei Costruttori, Modificatori e SimpleDate v3}
\author{\textbf{Docente:} Prof. Michele Pasqua \quad---\quad \textbf{Appunti Revisione:} Wally}
\date{A.A. 2025/2026 -- Universit\`a degli Studi di Verona}

\begin{document}
\maketitle
\thispagestyle{fancy}
\vspace{-1.2em}
\noindent\textcolor{brandblue}{\rule{\textwidth}{0.8pt}}
\vspace{1.2em}

In questo blocco analizziamo in dettaglio la teoria dei costruttori,
dell'incapsulamento e dell'information hiding
(\href{file:///home/wally/Documenti/GitHub/Programmazione-II/25-26/Teoria-e-Esercitazioni/Slides-20260902/T08\%20-\%20Class\%20Constructors\%20and\%20Encapsulation.pdf}{T08
- Class Constructors and Encapsulation.pdf}), confrontandola con il
codice ufficiale in
\href{file:///home/wally/Documenti/GitHub/Programmazione-II/25-26/Teoria-e-Esercitazioni/Codice-20260902/Lezione08/SimpleDate/src}{\passthrough{\lstinline!Lezione08!}},
dove nasce l'enum \textbf{\passthrough{\lstinline!Language.java!}},
vengono aggiunti gli array dei mesi per la stampa estesa
(\passthrough{\lstinline!prettyPrint()!}) e si introduce
l'\textbf{Enhanced For (for-each)}.

\begin{center}\rule{0.5\linewidth}{0.5pt}\end{center}

\subsubsection{1. Teoria: Concetti Teorici dalle Slide (Slide
T08)}\label{teoria-concetti-teorici-dalle-slide-slide-t08}

\paragraph{A. Il Ciclo di Vita e i Meccanismi dei Costruttori (Slide
1--12)}\label{a.-il-ciclo-di-vita-e-i-meccanismi-dei-costruttori-slide-112}

\begin{itemize}
\tightlist
\item
  \textbf{Cosa avviene esattamente dietro le quinte con
  \passthrough{\lstinline!new Clazz()!}}:

  \begin{enumerate}
  \def\labelenumi{\arabic{enumi}.}
  \tightlist
  \item
    La JVM calcola l'ingombro in byte di tutti i campi e alloca un
    blocco contiguo di memoria nello \textbf{Heap}.
  \item
    Genera un riferimento (\emph{reference}) univoco che punta a
    quell'area di memoria.
  \item
    Inizializza tutti i campi ai rispettivi valori di default
    (\passthrough{\lstinline!0!}, \passthrough{\lstinline!false!},
    \passthrough{\lstinline!null!}).
  \item
    Invoca il \textbf{costruttore} designato passando implicitamente il
    puntatore \passthrough{\lstinline!this!}.
  \item
    Restituisce l'indirizzo dell'oggetto alla variabile a sinistra
    dell'assegnamento.
  \end{enumerate}
\item
  \textbf{Caratteristiche Tassative di un Costruttore}:

  \begin{itemize}
  \tightlist
  \item
    Deve avere lo \textbf{stesso identico nome della classe}.
  \item
    \textbf{Non ha alcun tipo di ritorno} (nemmeno
    \passthrough{\lstinline!void!}!).
  \item
    Viene eseguito una sola volta, all'atto della creazione
    dell'istanza.
  \end{itemize}
\item
  \textbf{Default Constructor (Costruttore di Default)}:

  \begin{itemize}
  \tightlist
  \item
    Se il programmatore \textbf{non dichiara alcun costruttore}, il
    compilatore Java inserisce automaticamente un costruttore vuoto
    senza argomenti (\passthrough{\lstinline!public Clazz() \{ \}!}).
  \item
    Attenzione: \textbf{Regola aurea d'esame}: se il programmatore
    dichiara \emph{anche un solo} costruttore con parametri (es.
    \passthrough{\lstinline!public Car(String color)!}), il compilatore
    \textbf{non genera più il costruttore di default}! Invocare
    \passthrough{\lstinline!new Car()!} senza argomenti provocherà un
    \textbf{Compile-Time Error}.
  \end{itemize}
\item
  \textbf{False Constructor (Il Falso Costruttore --- Trabocchetto
  d'Esame)}:

  \begin{itemize}
  \item
    Se si inserisce un tipo di ritorno (ad esempio
    \passthrough{\lstinline!void Car(String color)!}):

\begin{lstlisting}[language=Java]
public class Car {
    void Car(String color) { ... } // Attenzione:  NON È UN COSTRUTTORE!
}
\end{lstlisting}
  \item
    Per il compilatore questo è un \textbf{normale metodo d'istanza}
    (avente incidentalmente lo stesso nome della classe). Il compilatore
    aggiungerà comunque il costruttore di default
    \passthrough{\lstinline!Car()!}, e il metodo
    \passthrough{\lstinline!void Car(...)!} non verrà mai invocato al
    momento del \passthrough{\lstinline!new!}!
  \end{itemize}
\item
  \textbf{Copy Constructor}:

  \begin{itemize}
  \tightlist
  \item
    Costruttore che accetta un'altra istanza della medesima classe per
    clonarne lo stato (\passthrough{\lstinline!Car(Car other)!}). In
    Java non esiste la copia automatica bit-a-bit del C++: deve essere
    scritta esplicitamente dal programmatore.
  \end{itemize}
\item
  \textbf{Constructor Overloading e Chaining con
  \passthrough{\lstinline!this(...)!}}:

  \begin{itemize}
  \tightlist
  \item
    Più costruttori con firme distinte.
  \item
    La chiamata \passthrough{\lstinline!this(...)!} verso un costruttore
    fratello serve a evitare duplicazione di logica di validazione e
    \textbf{deve essere tassativamente la prima riga di codice nel corpo
    del costruttore}.
  \item
    \emph{Meme delle slide (p.~11)}: Il matematico in lacrime che grida
    \emph{``Abuse of notation''} vs il Chad Programmer che usa
    l'overloading per semplificare il codice.
  \end{itemize}
\item
  \textbf{Distruttori e il Metodo \passthrough{\lstinline!finalize()!}}:

  \begin{itemize}
  \tightlist
  \item
    In Java non esiste la deallocazione manuale: non esistono
    distruttori (\passthrough{\lstinline!\~Car()!}).
  \item
    Esisteva il metodo protetto
    \passthrough{\lstinline!protected void finalize()!} invocato prima
    della rimozione da parte del Garbage Collector.
  \item
    Attenzione: \textbf{Attenzione}:
    \passthrough{\lstinline!finalize()!} è ufficialmente
    \textbf{deprecato da Java 9} in poi (e rimosso/reso no-op nelle
    versioni moderne) per l'imprevedibilità temporale dell'esecuzione
    del GC e rischi di deadlock. Non va mai utilizzato.
  \end{itemize}
\end{itemize}

\begin{center}\rule{0.5\linewidth}{0.5pt}\end{center}

\paragraph{B. Incapsulamento, Information Hiding e Modificatori di
Accesso (Slide
13--24)}\label{b.-incapsulamento-information-hiding-e-modificatori-di-accesso-slide-1324}

\begin{itemize}
\tightlist
\item
  \textbf{Incapsulamento}: aggregazione fisica nello stesso file
  sorgente di dati (campi) e procedure (metodi).
\item
  \textbf{Information Hiding}: occultamento della rappresentazione
  interna dello stato.

  \begin{itemize}
  \tightlist
  \item
    \emph{Meme della Cipolla (p.~18)}: \emph{«Most software is written
    like an onion: The more layers you peel back, the more you want to
    cry»}. L'Information Hiding crea confini netti per proteggere il
    codice dal disfacimento.
  \item
    \emph{Meme di Gandalf (p.~24)}: \emph{«YOU SHALL NOT ACCESS MY
    MEMBERS»}.
  \end{itemize}
\item
  \textbf{I 4 Livelli di Visibilità in Java}:

  \begin{enumerate}
  \def\labelenumi{\arabic{enumi}.}
  \tightlist
  \item
    \passthrough{\lstinline!private!}: visibile esclusivamente
    all'interno della stessa classe.
  \item
    \emph{Default (Package-Private)}: nessun modificatore; visibile
    all'interno della classe e a tutte le classi residenti nello stesso
    package.
  \item
    \passthrough{\lstinline!protected!}: visibile nel package e a tutte
    le sottoclassi derivate (anche in package diversi).
  \item
    \passthrough{\lstinline!public!}: visibile ovunque, da qualsiasi
    package del Classpath.
  \end{enumerate}
\item
  \textbf{Vantaggi Architetturali di Getter e Setter}:

  \begin{itemize}
  \tightlist
  \item
    \emph{Nei Setter}: possibilità di intercettare valori anomali,
    applicare regole di business e loggare le modifiche prima di mutare
    lo stato.
  \item
    \emph{Nei Getter}: possibilità di alterare la rappresentazione
    interna della memoria senza rompere il codice dei client
    (\textbf{Retro-compatibilità / Loose Coupling}). Ad esempio,
    cambiare il campo interno da \passthrough{\lstinline!int age!} a
    \passthrough{\lstinline!short age!} per dimezzare l'occupazione di
    RAM, mantenendo il getter
    \passthrough{\lstinline!public int getAge() \{ return age; \}!}
    (promozione implicita di tipo) senza che il mondo esterno debba
    ricompilare o cambiare una sola riga di codice.
  \end{itemize}
\end{itemize}

\begin{center}\rule{0.5\linewidth}{0.5pt}\end{center}

\paragraph{C. Asserzioni e Testing (Slide
25--26)}\label{c.-asserzioni-e-testing-slide-2526}

\begin{itemize}
\item
  La parola chiave \passthrough{\lstinline!assert!}:

\begin{lstlisting}[language=Java]
assert person.getAge() != 0 : "Age cannot be 0";
\end{lstlisting}

  \begin{itemize}
  \item
    Se l'espressione a sinistra valuta \passthrough{\lstinline!false!},
    la JVM interrompe l'esecuzione e lancia un
    \passthrough{\lstinline!AssertionError!} mostrando il messaggio a
    destra dei due punti.
  \item
    Richiede obbligatoriamente l'attivazione a runtime con il flag JVM
    \textbf{\passthrough{\lstinline!-ea!}} (\emph{Enable Assertions}):

\begin{lstlisting}[language=bash]
java -ea TestPerson
\end{lstlisting}
  \end{itemize}
\end{itemize}

\begin{center}\rule{0.5\linewidth}{0.5pt}\end{center}

\subsubsection{\texorpdfstring{2. Codice: Evoluzione del Codice:
\href{file:///home/wally/Documenti/GitHub/Programmazione-II/25-26/Teoria-e-Esercitazioni/Codice-20260902/Lezione08/SimpleDate/src}{\texttt{Lezione08}}}{2. Codice: Evoluzione del Codice: Lezione08}}\label{codice-evoluzione-del-codice-lezione08}

In \passthrough{\lstinline!Lezione08!} il docente applica un refactoring
strutturale a \passthrough{\lstinline!SimpleDate!}:

\paragraph{\texorpdfstring{1. Introduzione del Tipo Enumerato
\href{file:///home/wally/Documenti/GitHub/Programmazione-II/25-26/Teoria-e-Esercitazioni/Codice-20260902/Lezione08/SimpleDate/src/Language.java}{\texttt{Language.java}}}{1. Introduzione del Tipo Enumerato Language.java}}\label{introduzione-del-tipo-enumerato-language.java}

Viene abbandonato il primitivo \passthrough{\lstinline!byte lang!} (dove
\passthrough{\lstinline!0!} era IT e \passthrough{\lstinline!1!} era US,
fragile e poco leggibile) a favore di un tipo enum fortemente tipizzato:

\begin{lstlisting}[language=Java]
public enum Language {
    IT, US;

    public String toString() {
        switch (this) {
            case IT: return "italiano";
            case US: return "american";
            default: return null;
        }
    }
}
\end{lstlisting}

\emph{Finezza}: Anche gli enum in Java sono tipi reference completi e
possono sovrascrivere \passthrough{\lstinline!toString()!} per
restituire una descrizione testuale appropriata.

\begin{center}\rule{0.5\linewidth}{0.5pt}\end{center}

\paragraph{\texorpdfstring{2. Refactoring di
\href{file:///home/wally/Documenti/GitHub/Programmazione-II/25-26/Teoria-e-Esercitazioni/Codice-20260902/Lezione08/SimpleDate/src/Date.java}{\texttt{Date.java}}}{2. Refactoring di Date.java}}\label{refactoring-di-date.java}

\begin{lstlisting}
 public class Date {
-    private byte lang; // 0: IT, 1: US
+    private Language lang;
+    private static final String[] MONTHS_IT = { "gennaio", "febbraio", "marzo", "aprile", "maggio", "giugno", "luglio", "agosto", "setembre", "ottobre", "novembre", "dicembre" };
+    private static final String[] MONTHS_US = { "January", "February", "March", "April", "May", "June", "July", "August", "September", "October", "November", "December"};

     public Date(int day, int month, int year) {
         this.day = day;
         this.month = month;
         this.year = year;
         verify();
-        lang = 0;
+        lang = Language.IT;
     }

-    public void setLang(byte lang) { ... }
+    public void setAmerican() { lang = Language.US; }
+    public void setItalian() { lang = Language.IT; }

+    public String getMonthAsString() {
+        if (lang == Language.IT) return MONTHS_IT[month-1];
+        else return MONTHS_US[month-1];
+    }

+    public String prettyPrint() {
+        if (lang == Language.IT) return day + " " + MONTHS_IT[month-1] + " " + year;
+        else return MONTHS_US[month-1] + " " + day + ", " + year;
+    }
\end{lstlisting}

\paragraph{\texorpdfstring{Analisi: Dettagli Tecnici e Trabocchetti nel
Codice di
\texttt{Date}:}{Analisi: Dettagli Tecnici e Trabocchetti nel Codice di Date:}}\label{analisi-dettagli-tecnici-e-trabocchetti-nel-codice-di-date}

\begin{enumerate}
\def\labelenumi{\arabic{enumi}.}
\item
  \textbf{Separazione Semantica dei Metodi di Configurazione}: Invece di
  un setter generico che accetta valori arbitrari, il prof definisce due
  metodi espliciti auto-esplicativi:
  \passthrough{\lstinline!setItalian()!} e
  \passthrough{\lstinline!setAmerican()!}.
\item
  \textbf{Array di Costanti di Classe
  (\passthrough{\lstinline!private static final String[]!})}: Gli
  elenchi dei nomi dei mesi sono \passthrough{\lstinline!static final!}
  (condivisi da tutte le istanze senza duplicare memoria nello Heap) e
  \passthrough{\lstinline!private!} (inaccessibili all'esterno).
  \emph{Curiosità nel codice del prof}: nell'array
  \passthrough{\lstinline!MONTHS\_IT!} c'è un piccolo refuso:
  \passthrough{\lstinline!"setembre"!} con una sola `t'.
\item
  \textbf{Persistenza del Refuso del Prof nei Getter}: Anche in questa
  lezione, alle righe 43--44 di \passthrough{\lstinline!Date.java!}:

\begin{lstlisting}[language=Java]
public int getDay() { return day; }
public int getMonth() { return day; } // Ritorna day anziché month!
public int getYear() { return day; }  // Ritorna day anziché year!
\end{lstlisting}

  Un chiaro refuso di copia-incolla che non è stato corretto dal docente
  nel rilascio ufficiale. Nel proprio codice personale è doveroso
  restituire i rispettivi campi \passthrough{\lstinline!month!} e
  \passthrough{\lstinline!year!}.
\item
  \textbf{\passthrough{\lstinline!prettyPrint()!} e Formattazione
  Culturale}:

  \begin{itemize}
  \tightlist
  \item
    In italiano: giorno, mese per esteso, anno (es.
    \passthrough{\lstinline!7 gennaio 2025!}).
  \item
    In inglese/americano: mese per esteso, giorno, virgola, anno (es.
    \passthrough{\lstinline!February 20, 2025!}).
  \end{itemize}
\end{enumerate}

\begin{center}\rule{0.5\linewidth}{0.5pt}\end{center}

\paragraph{\texorpdfstring{3. Novità in
\href{file:///home/wally/Documenti/GitHub/Programmazione-II/25-26/Teoria-e-Esercitazioni/Codice-20260902/Lezione08/SimpleDate/src/MainDate.java}{\texttt{MainDate.java}}:
Array di Oggetti ed Enhanced
For}{3. Novità in MainDate.java: Array di Oggetti ed Enhanced For}}\label{novituxe0-in-maindate.java-array-di-oggetti-ed-enhanced-for}

\begin{lstlisting}[language=Java]
// Inizializzazione rapida di un array di oggetti nello Heap
Date[] dates = { d1, d2, d3, new Date(14, 2, 2024) };

// 1. Scansione all'indietro classica tramite indice
for (int i = dates.length - 1; i >= 0; i--)
    System.out.println(dates[i].toString() + ": " + dates[i].getMonthAsString());

// 2. Scansione in avanti con Enhanced For (For-Each)
for (Date date : dates)
    System.out.println(date.toString() + ": " + date.getMonthAsString());
\end{lstlisting}

\begin{itemize}
\tightlist
\item
  Il costrutto
  \textbf{\passthrough{\lstinline!for (Tipo elemento : collezione)!}}
  evita di manipolare manualmente indici di scorrimento, rendendo il
  codice più leggibile ed eliminando il rischio di errori di off-by-one
  o out-of-bounds.
\end{itemize}

\begin{center}\rule{0.5\linewidth}{0.5pt}\end{center}

\subsubsection{3. Ciclo: Corrispondenza con i Tuoi Moduli
Workspace}\label{ciclo-corrispondenza-con-i-tuoi-moduli-workspace}

Nel tuo repository
\href{file:///home/wally/Documenti/GitHub/Programmazione-II/25-26/esercizi-wally-25-26}{\passthrough{\lstinline!esercizi-wally-25-26!}}:
- \textbf{\passthrough{\lstinline!es06!}}:
\href{file:///home/wally/Documenti/GitHub/Programmazione-II/25-26/esercizi-wally-25-26/es06/src/main/java/es06/Person.java}{\passthrough{\lstinline!Person.java!}}
implementa l'esempio accademico della persona con i metodi di
validazione \passthrough{\lstinline!verifyAge!} e
\passthrough{\lstinline!verifyName!}, incapsulamento rigoroso e gestione
dei setter/getter. - \textbf{\passthrough{\lstinline!es07!}}:
\href{file:///home/wally/Documenti/GitHub/Programmazione-II/25-26/esercizi-wally-25-26/es07/src/main/java/es07/Language.java}{\passthrough{\lstinline!Language.java!}}
definisce l'enum \passthrough{\lstinline!Language!} sfruttando
modernamente la \emph{Switch Expression} di Java 14+
(\passthrough{\lstinline!case IT -> "italiano";!}).

\begin{center}\rule{0.5\linewidth}{0.5pt}\end{center}

\begin{studynote}[Nota di Studio] Con la \textbf{Lezione 08} abbiamo completato l'intero ciclo
di vita dei costruttori, l'incapsulamento dei dati, gli enum base e la
gestione dei formati estesi.

Il prossimo blocco è la \textbf{Lezione 09 / Slide T09: \emph{Arrays and
Enumerative Types}}: - Approfondimento teorico sistematico sugli
\textbf{Array in Java} (dichiarazione, allocazione nello Heap, proprietà
immutabile \passthrough{\lstinline!.length!}, array multidimensionali,
copie superficiali vs profonde). - Analisi rigorosa degli
\textbf{Enumerative Types (\passthrough{\lstinline!enum!})}: costruttori
di enum, campi interni, metodi \passthrough{\lstinline!values()!} e
\passthrough{\lstinline!ordinal()!}. - Evoluzione del codice in
\passthrough{\lstinline!Lezione09!}: - Potenziamento di
\passthrough{\lstinline!Language.java!} con campi interni
(\passthrough{\lstinline!format!},
\passthrough{\lstinline!description!}) e costruttore privato. -
Riscrittura ed espansione di \passthrough{\lstinline!Date.java!}.
\end{studynote}

\textbf{Dammi conferma per aprire la Lezione 09 / T09 e continuare!}


\end{document}
